        	\subsection{Geometry of half-bubbles}\label{subs:geomhalf}
        	
        	Let us discuss now some basic facts about the geometry and topology of half-bubbles. First of all, keeping in mind the quantization identity presented in Theorem 1, we wish to express the total curvature of a half-bubble in terms of its topology.
        	Clearly, in ambient dimension three ($n=2$) we can apply the Gauss equations to get
        	\begin{equation}\label{eq:half2nd}
        	2{\jepPubScript{A}}(\Sigma)={\jepPubScript{A}}(\check{\Sigma})=-2\int_{\check{\Sigma}} K_{\check{\Sigma}}\,d\h^2,
        	\end{equation}
        	while by Gauss-Bonnet (cf. Jorge-Meeks \cite{JM83}), denoting by $\gamma(\check{\Sigma})$ the genus of $\check{\Sigma}$ and by $b(\check{\Sigma})$ the number of its ends, we have
        	\begin{equation}\label{eq:GBcomplete}
        	\int_{\check{\Sigma}} K_{\check{\Sigma}}\,d\h^2=2\pi (\chi(\check{\Sigma})-b(\check{\Sigma}))=2\pi(2-2\gamma(\check{\Sigma})-2b(\check{\Sigma})) .
        	\end{equation}
        	It is then important to relate such data to the topological data of the half-bubble $\Sigma$. There are two easy examples of half-bubbles to be kept in mind throughout the following discussion, which we will prove in Lemma \ref{lem:euler} to be, roughly speaking, prototypical.
        	
        	\begin{definition}\label{def:horcurvercut} In the flat Euclidean space $\R^3$ we introduce the following terminology:
        		\begin{itemize} \item{the half-catenoid given by
        			\[
        			\left\{(x^1, x^2, x^3)\in \R^3 : \ (x^3)^2=\cosh ((x^1)^2+(x^2)^2), \ \ x^3\geqslant 0 \right\}
        			\]
        			will be called \textsl{horizontally cut half-catenoid};}	
       \item{the half-catenoid given by
        			\[
        		\left\{	(x^1, x^2, x^3)\in \R^3 : \ (x^3)^2=\cosh ((x^1)^2+(x^2)^2), \ \ x^1\geqslant 0 \right\}	
        			\]
        			will be called \textsl{vertically cut half-catenoid}.}
        			\end{itemize}
        	\end{definition}
        	
        	 Before proceeding further, it is helpful to recall some information about the asymptotic structure of complete minimal surfaces in $\R^3$.
        	
        	\begin{rmk}\label{rem:reginf}
        		Let $S$ be a complete, embedded minimal surface of $\R^3$. Then:
        		\begin{itemize}
        			\item{$S$ has finite total curvature if and only if it has finite Morse index (cf. \cite{FC85});}
        			\item{if $S$ satisfies either of these two equivalent assumptions then it is regular at infinity (cf. \cite{Sc83}), namely it can be decomposed, outside any sufficiently large compact set, into a finite number of connected components (named \textsl{ends}) and each such end can be described as a graph (over a plane in $\R^3$) of a defining function having an expansion of the form
        				\[
        				u(x')=a\log|x'|+b+\frac{c_1 x^1}{|x|^2}+\frac{c_2 x^2}{|x|^2}+O(|x'|^{-2}), \ |x'|\longrightarrow\infty
        				\]
        				in suitable coordinates $x=(x^1,x^2,x^3)$, where $x'=(x^1,x^2)$.
        			}
        		\end{itemize}
        		Notice that if $\Sigma$ is a half-bubble then both facts apply to its double $\check{\Sigma}$ and hence, a posteriori, one can talk about the ends of a half-bubble and this notion is well-defined. Furthermore, each end of a half-bubble will have a precise asymptotic description of the type above.
        	\end{rmk}
        	
        		\begin{rmk}
        			We can somewhat strengthen Lemma \ref{lem:II1} and extend to half-bubbles the aforementioned equivalence result by Fischer-Colbrie:
        			\end{rmk}
        		\begin{prop}\label{pro:equiv}
        			A two-dimensional half-bubble has finite index if and only if it has finite total curvature.	The same conclusion holds true in all dimensions provided Euclidean volume growth is assumed.
        		\end{prop}
        		One of the two implications has been proved above in Lemma \ref{lem:II1}, while the other follows by combining Lemma \ref{lem:specadd} and the inequalities in \eqref{eq:fundineq}.
        	
        	\
        	
        	Let us now proceed in our discussion, relating all half-bubbles to one of the two models given in Definition \ref{def:horcurvercut}.
        	
        	\begin{lem}\label{lem:euler}
        		Let $\Sigma$ be a half-bubble and let $\check{\Sigma}$ denote its double in $\R^3$. Then one of the following two alternative cases occurs:
        		\begin{enumerate}
        			\item{the number of ends of $\check{\Sigma}$ equals double the number of ends of $\Sigma$, and their Euler characteristics are related by the equation
        				\[
        				\chi(\check{\Sigma})=2\chi(\Sigma);
        				\]	
        			}
        			\item{the number of ends of $\check{\Sigma}$ equals the number of ends of $\Sigma$,  and their Euler characteristics are related by the equation
        				\[
        				\chi(\check{\Sigma})=2\chi(\Sigma)-b(\Sigma).
        				\]
        			}	
        			\end{enumerate}
        			In either case, one has that $\chi(\check{\Sigma})-b(\check{\Sigma})=2(\chi(\Sigma)-b(\Sigma))$.
        	\end{lem}	
        	
        	Clearly, both cases do occur: (1) is exemplified by the horizontally cut half-catenoid, while (2) is exemplified by the vertically cut half-catenoid. 
        	
        	\begin{proof}
        		Based on Remark \ref{rem:reginf} applied to the symmetrized bubble $\check{\Sigma}$ we have that either each of its ends is contained in one of the half-spaces $\Pi_1$ and $\Pi_2$ (provided we remove from $\check{\Sigma}$ a sufficiently large ball centered at the origin) or instead all of them intersect $\Pi$. By symmetry, it is clear that in the first case the number of ends of $\check{\Sigma}$ equals double the number of ends of $\Sigma$ and we can justify the equation for the Euler characteristic using the Gauss-Bonnet theorem: since $\Sigma$ is assumed to have finite total curvature we can write for $\Sigma_r:=\Sigma \cap \left\{|x|\leqslant r\right\}$
        		\[
        		\int_{\Sigma}K_{\Sigma}\,d\h^2=\lim_{r\to\infty}\int_{\Sigma_r} K_{\Sigma}\,d\h^2
        		\]
        		and we can compute the limit on the right-hand side along a sequence of generic radii $\left\{r_i\right\}$, so that the intersection $\Sigma\cap \left\{|x|=r_i\right\}$ is transverse and thus consisting of finitely many, smooth closed curves, hence
        		\[
        		\int_{\Sigma_{r_i}} K_{\Sigma}\,d\h^2=2\pi \chi(\Sigma)-2\pi b(\Sigma)+o(1)
        		\]
        		(where we have used the fact that the geodesic curvature of $\partial\Sigma\subset\Sigma$ is zero at all points).
        		Thus the conclusion comes straight from combining this equation with \eqref{eq:half2nd} and \eqref{eq:GBcomplete}. This is the scenario corresponding to case (1).
        		
        		If instead all ends of $\check{\Sigma}$ intersect the plane $\Pi$ then we are in case (2) and the claim concerning $\chi(\check{\Sigma})$ comes, once again, by an elementary computation: this time we have, along a sequence of generic radii 
        		\[
        		\int_{\Sigma_{r_i}} K_{\Sigma}\,d\h^2=2\pi \chi(\Sigma)-\pi b(\Sigma)-\pi b(\Sigma)+o(1)
        		\]
        		by keeping the contributions of both the geodesic curvature and of the exterior angles (at the non-smooth points of $\Sigma_r$) into account. In that respect, notice that  there are exactly $2b(\Sigma)$ exterior angles, each of them being $\pi/2+o(1)$ as $i\to\infty$. Thereby, the conclusion follows.
        	\end{proof}	
        	
        	
        	
        		
        		
        		
        		
        	
        
        	
        	
        	We can extend to half-bubbles some further classification results that are well-known for complete minimal surfaces in $\R^3$. This discussion parallels the one presented in the previous subsection, which was based on index-theoretic criteria instead.
        	
        	
        	\begin{cor}\label{cor:oneend}
        		A two-dimensional half-bubble with one end is isometric to either a half-plane or to a horizontally cut half-catenoid.
        	\end{cor}	
        	
        	\begin{proof}
        		Denoted by $\Sigma$ the half-bubble in question and $\check{\Sigma}$ its double, Lemma \ref{lem:euler} implies that \textsl{either} $\Sigma$ and $\check{\Sigma}$ have the same number of ends (one) \textsl{or} $\check{\Sigma}$ will have two ends. By classical results of Schoen (see \cite{Sc83}) the first case happens if and only if $\check{\Sigma}$ is a plane, and the second if $\check{\Sigma}$ is a catenoid. The conclusion in the former case is straightforward, so let us discuss the latter. The only way $\Sigma$ has exactly one end (which we are assuming) is that it is cut with a plane that is parallel to its two ends. But then the only way such a cut gives rise to a free boundary minimal surface is when the plane in question passes through the center of the catenoid, which proves the claim.
        	\end{proof}	
        	
        	\begin{cor}\label{cor:genuszero}
        		A two-dimensional half-bubble whose Euler characteristic equals one is isometric to either a half-plane or to a vertically cut half-catenoid.
        	\end{cor}	
        	
        	\begin{proof}
        		Like we did above, we need to consider two cases and apply the equations provided by Lemma \ref{lem:euler} for the Euler characteristics. If we are in case (1), the symmetrized bubble $\check{\Sigma}$ satisfies $\chi(\check{\Sigma})=2$ and this is not possible because the standard Jorge-Meeks formula for the Euler characteristic (cf. the second equality in \eqref{eq:GBcomplete}) forces $\chi(\check{\Sigma})\leqslant 1$. 
        		Instead, in case (2) we would get that the bubble $\check{\Sigma}$ has genus zero, which allows to invoke \cite{LR91}: $\check{\Sigma}$ is either a plane or a catenoid. At this stage, the conclusion comes at once by considering all possible planes of symmetry of a catenoid.
        	\end{proof}	
        	
        	
        	 	
        	 	
